% Generated draft: requires mathematical and source review.
\documentclass[11pt]{article}
\usepackage[margin=25mm]{geometry}
\usepackage{fontspec}
\setmainfont{Latin Modern Roman}
\usepackage{amsmath,amssymb,mathtools,unicode-math}
\setmathfont{Latin Modern Math}
\usepackage{enumitem,needspace,xurl,hyperref}
\setlength{\parindent}{0pt}
\setlength{\parskip}{5pt}
\setlength{\emergencystretch}{3em}
\newcommand{\N}{\mathbb{N}}
\newcommand{\Z}{\mathbb{Z}}
\newcommand{\Q}{\mathbb{Q}}
\newcommand{\R}{\mathbb{R}}
\newcommand{\C}{\mathbb{C}}
\newcommand{\F}{\mathbb{F}}
\newcommand{\A}{\mathbb{A}}
\newcommand{\PP}{\mathbb{P}}
\newcommand{\E}{\mathbb{E}}
\newcommand{\eps}{\varepsilon}
\newcommand{\dd}{\,\mathrm{d}}
\begin{document}
\section*{UnsolvedMath ID: 20000320}
\textbf{Source problem code:} catalogue entry
\subsection{Definitions and mathematical statement}
For a surface family \(f:X\to B\), let \(\omega_{X/B}^{[m]}=(\omega_{X/B}^{\otimes m})^{**}\) denote the \(m\)-th reflexive power. The original defines base-change compatibility by requiring, for every \(\phi:B'\to B\), with \(X'=X\times_BB'\) and induced \(\phi_X:X'\to X\),
\[
(\phi_X^*\omega_{X/B})^{[m]}\cong\phi_X^*(\omega_{X/B}^{[m]}).
\]
Koll\'ar's condition requires this for every integer \(m>0\); where the relative dualizing sheaf itself has the appropriate base-change identification, the usual expression is compatibility of \(\omega_{X/B}^{[m]}\) with \(\omega_{X'/B'}^{[m]}\). Relative \(\Q\)-Cartierness means a suitable positive reflexive power is locally invertible. The AIM question asks whether relative \(\Q\)-Cartierness is equivalent to Koll\'ar's condition. The general family category and the relative-reflexive conventions on arbitrary test bases need to be specified.

The original says the equivalence is known for smoothings over a discrete valuation ring, reference [Hac], Proposition 10.14, and that it fails in positive characteristic, as Koll\'ar showed at the conference. It attributes the question implicitly to Koll\'ar, contrasting his stronger condition in [Kol90] with the weaker moduli-functor definition used with Shepherd-Barron in [KSB88]. These are source-specific historical qualifications.

The attached report claims that in characteristic zero, for a flat finite-type family of semi-log-canonical surfaces over a \emph{reduced original base}, the two conditions are equivalent. It says Koll\'ar's hull criterion gives compatibility with every test base change, including nonreduced tests. Reducedness is imposed on \(B\), not on every allowed \(B'\). It separately invokes Lee--Nakayama examples to say reducedness alone is insufficient without the slc hypothesis. The complete hull and family hypotheses are not supplied.

For a nonreduced original base, the report gives the Altmann--Koll\'ar first-order deformation associated to the interval \([-2/5,2/5]\), in Section 5.7 of their referenced work. Over the first-order base \(\operatorname{Spec}k[\varepsilon]/(\varepsilon^2)\), its asserted behavior is
\[
\omega_{X/B}^{[5]}\text{ invertible},\qquad
\omega_{X/B}^{[3]}\text{ fails base change to the closed fiber}.
\]
The source calls this deformation ``VW-but-not-qG'' and identifies its singularity as \(\tfrac1{20}(1,11)\). These deformation labels must not be conflated with the weaker local invertible-power condition being tested.

Its candidate lattice witness is that \((1,2/5)\) lies in the relevant half-open zone and the exponent-three coset \(M-(3/5)\overline R\), but not on the permitted diagonal. The report says this violates condition \((*)_3\), hence third-power base change, while the fifth power remains invertible. The lattice \(M\), vector \(\overline R\), zone, diagonal, and condition \((*)_3\) are not defined in the supplied summary; the numerical witness is preserved without pretending these data form a complete explicit counterexample here.
\subsection{Short English statement}
Compare relative invertibility of one canonical power with base-change compatibility of every positive reflexive power. The report claims equivalence for characteristic-zero slc families over reduced bases, but gives a nonreduced first-order example with invertible fifth power and failing third-power base change.
\subsection{Sources}
\begin{itemize}
\item AIM workshop ``Compact moduli spaces and birational geometry,'' Problem 1.1: \url{https://aimath.org/WWN/birational/birational.pdf}. No PDF was accessed here.
\item Local canonical registry \texttt{aim-algebraic-geometry-notes.json}, \texttt{notes[319]}, tag \texttt{problem}; permanent UnsolvedMath ID \(20000320\), source code \texttt{catalogue entry}. Original references are [Hac], Proposition 10.14, [Kol90], and [KSB88], with the historical Koll\'ar/Shepherd-Barron attribution and positive-characteristic conference counterexample; their full bibliography is not supplied. The verification field claims displaced OCR subscripts and superscripts can be recovered. Original remarks and literature field are otherwise empty.
\item Supplied research attempt 1 refers to Koll\'ar's hull criterion, Lee--Nakayama examples, and Altmann--Koll\'ar Section 5.7, and proposes the exponent-three lattice witness. No titles, identifiers, or full authoritative references are supplied. The report is machine generated, not peer reviewed, and requires independent verification before citation.
\end{itemize}
\end{document}
